\subsection{The ordinary-node defect in transverse type \texorpdfstring{$A_3$}{A3}}

\begin{statement}{P4.41}{Proved}
At a point of an actual support in Part~4.2, suppose the
entire reduced critical curve is the ordinary node
$C=\{z=xy=0\}$ and both branches have generic transverse
type $A_3$. Let $d_i,b_i,q$ be the Hessian zero orders,
projected conormal colengths, and normalization descent
defect of P4R.10. Then
\[
 \delta(C)+\sum_{i=1}^2\left(\frac{d_i}{2}+b_i\right)+q
                         \ge\frac92.                 \tag{4.41a}
\]
Here $R=0$ and $\delta(C)=1$. Equality makes the normalization
of the hypersurface smooth over the point.
\end{statement}
\begin{proof}
Put $P=(z,xy)$, and let $I$ be the ambient ideal of the
scheme-theoretic conductor image. We have $P^2=P^{(2)}\subset I$,
and $J=P/I$ is torsionfree of rank one on both branches.
After dividing its normalized image by the two branchwise
vanishing orders, it is an $A=\mathcal O_C$ submodule of
$B=\mathbb C\{x\}\oplus\mathbb C\{y\}$ generating both
summands. The conductor of $A$ in $B$ is $\mathfrak m_A$,
so $\mathfrak m_AB\subset J$. Thus $J$ is determined by
a subspace of $\mathbb C^2$ projecting onto both coordinates.
If that subspace is a line, rescaling gives $J\simeq A$ and
$q=1$. If it is the whole space, $J\simeq B$ and $q=0$.
These are actual module isomorphisms. The divided orders are
the $b_i$ already in the formula.

We use the ring condition of the actual conductor:
\[
 \operatorname{Hom}_{\mathcal O_S}(I,I)=
 \operatorname{Hom}_{\mathcal O_S}(I,\mathcal O_S).
                                                               \tag{4.41b}
\]
The right side is the normalization algebra by finite duality,
and $I$ is an ideal of that algebra. For the determinant
presentation of a Cohen--Macaulay codimension-two conductor
inside a Gorenstein hypersurface, this condition is equivalent
to membership of the adjugate entries in $I$. This is precisely
de Jong--van Straten, Theorem~1.12, printed p.~532
\cite{DeJongVanStraten1990Normalization}. Its hypotheses hold
here: the ambient space is regular, the reduced support is a
hypersurface, and its conductor quotient is Cohen--Macaulay.
We give all the matrices used below.

Suppose first that $q=1$. The surjection
$P/P^2\to J$ is a surjection $A^2\to A$, with primitive
kernel generator $g$. Taking a complementary generator $k$
of $P$ gives $I=(g,k^2)$. A determinant presentation
$f=ag+bk^2$ has two-by-two adjugate containing $a,b,g,k^2$.
Condition (4.41b) says $a,b\in I$, so the exact equation is
\[
                   f=Ag^2+Bgk^2+Ck^4.                \tag{4.41c}
\]
We claim that $A(0)$ cannot be a unit. If $g_z(0)\ne0$,
straighten $g$, take $k$ to be a nodal equation in $g=0$,
and apply the parameter Morse lemma in $g$. The residual
is $k^4$ times a unit after a small change in its fourth-order
coefficient. A morsification of the node and a shift of its
critical value away from zero give an isolated Morse point
of nonzero critical value.

If $g_z(0)=0$, its $xy$ coefficient is a unit and we take
$k=z$. Completing the square in (4.41c), including the
square root of $A$ in its first factor, gives exactly
\[
                         f=G^2+z^4U,
\]
with $G(x,y,0)=xy$ times a unit and $G_z(0)=0$.
Add a small constant to $U$, then take the actual coordinate
$Z=zU^{1/4}$. The parameter Morse lemma for $G$ in $x,y$
gives $G=XY+a(Z)$, where $a\in(Z^2)$. A small change
$a\mapsto a+\eta Z^2$ makes $a''(0)\ne0$. Replacing
$G$ by $G-\epsilon$ produces at the origin a critical point
of value $\epsilon^2$. The $XY$ Hessian entry is
$-2\epsilon$, and the $ZZ$ entry is
$-2\epsilon a''(0)$, so the point is Morse.
These deformations retain type $A_3$ along the moving
critical curve at each fixed boundary. The nested boundary
transport of P4.28 and P4R.14 places its square-$-2$ sphere
in the actual nearby torus, a contradiction. Hence $A(0)=0$.

If $g_z(0)$ is a unit, both saturated Hessian lines approach
$dz$, so $b_1=b_2=1$ and $d_i\ge1$. If $g_z(0)=0$,
the differential of $g$ on each branch has exactly the
common factor $t_i$ and a remaining nonzero coefficient
in the other normal direction. Thus $b_i=0,d_i\ge3$.
In either case the left side of (4.41a) is at least five,
including $q=1$.

Now let $q=0$. Choose actual generators $a,b$ of $P$ mapping
to $(1,0),(0,1)$ in $J\simeq B$. Their coefficient matrix
is invertible modulo the maximal ideal, hence invertible
locally. We obtain
\[
 I/P^2=(ya,xb),\qquad
 f=xU(x)b^2+yV(y)a^2\pmod {P^3}.                      \tag{4.41d}
\]
The second identity records the Hessian kernel $b$ on the
$x$ branch and $a$ on the $y$ branch. Write the exact
relation $xy=\alpha a+\beta b$. Near a general point of
the $x$ axis, solve this for $y$ using coordinates $x,a,b$.
The null cubic coefficient contains
$V(0)\alpha(x,0)/x$; the remaining contribution from
$P^3$ is holomorphic at the node. Its identical vanishing
in type $A_3$ forces $V(0)\alpha(0)=0$. The other branch
gives $U(0)\beta(0)=0$.

If both $a_z(0),b_z(0)$ are nonzero, the relation for $xy$
has $\alpha(0)\beta(0)\ne0$. Hence $U(0)=V(0)=0$,
both $b_i=1$, and both $d_i\ge2$. The defect is at least
five. Otherwise interchange the branches so that $a_z(0)$
is a unit and $b_z(0)=0$. Then $\alpha(0)=0$ and
$\beta(0)$ is a unit, forcing $U(0)=0$. On the $x$ axis,
$db$ has exactly the common factor $x$ and remaining
$dy$ coefficient a unit. On the $y$ axis, $da$ is primitive
with a unit $dz$ coefficient. Thus
\[
 b_x=0,\quad d_x=3+\operatorname{ord}_xU\ge4,
 \qquad
 b_y=1,\quad d_y=1+\operatorname{ord}_yV\ge1.
\]
This proves (4.41a). Equality requires
$\operatorname{ord}_xU=1$ and $V(0)\ne0$.

We identify the normalization in this exact equality case.
Use $a$ as the ambient coordinate $z$. The generator relation
with $\beta$ a unit gives
\[
                    I=(a^2,ya,x^2y-x\alpha a).
\]
Modulo its first two generators, replace $\alpha$ by
$\alpha(x,0,0)=x\ell(x)$. The actual coordinate
$Y=y-\ell(x)a$ puts the ideal into
$I=(z^2,Yz,x^2Y)$. Renaming $Y$ as $y$, its Hilbert--Burch
matrix and determinant extension are
\[
 \begin{pmatrix}y&0\\-z&x^2\\0&-z\end{pmatrix},
 \qquad
 M=\begin{pmatrix}y&0&a\\-z&x^2&b\\0&-z&c\end{pmatrix},
 \qquad
 \det M=az^2+byz+cx^2y.
\]
The new two-by-two minors are
$yb+za,yc,-zc,-x^2a,-za,x^2c+zb$, up to signs.
Their membership in $I$ gives exactly
\[
             a=yA,\qquad c=zC,\qquad b=-x^2C+zB
                          \pmod I.
\]
Therefore $f\in(yz^2)+I^2$, so
\[
           f=Uyz^2+Vz^4+Wx^2y^2z+Tx^4y^2.            \tag{4.41e}
\]
The equality orders $d_y=1,d_x=4$ say precisely
$U(0)T(0)\ne0$.

Substitute $y=z^2v$ and divide (4.41e) by $z^4$.
At the finite root its $v$ derivative is $U(0)$.
The implicit-function theorem supplies the actual smooth
factor $g_1=y-z^2v(x,z)$. Division by it gives
\[
                  g_2=z^2A+x^2yzB+x^4yC,
                  \qquad A(0)C(0)\ne0.
\]
For clarity, this division retains the factors $x^2,x^4$:
subtracting the expression at $y=z^2v$ supplies one factor
$y-z^2v$; every extra substituted $y$ supplies $z^2$ and
is absorbed into the first term. The parameter Morse
centre in $z$ is divisible by $x^2y$. Substitution leaves
$x^4y$ times a unit. An actual coordinate change therefore
gives $g_2=Z^2+x^4Y$. Its finite birational normalization is
\[
                 (x,u)\longmapsto(x,-u^2,x^2u),
\]
with target coordinates $(x,Y,Z)$. It is smooth; finiteness
follows from $u^2=-Y$ and birationality from $u=Z/x^2$
away from $x=0$. The other branch $g_1$ is smooth already.
Thus the full normalization is smooth over the node.
\end{proof}

All determinant minors, the four minimum-defect cases,
the factor substitution, and the normalization equation
are verified by \texttt{code/part4/a3\_node.py}.
The output reports PASS. Its initial failed assertion
compared an expanded polynomial with a factored expression;
that output is retained, and the corrected assertion tests
their exact zero difference.

\begin{statement}{P4.42}{Proved}
In the requested global support classes, a positive-genus
ruled-resolution support whose critical curve has generic
type $A_3$ everywhere cannot have two two-branch points as
its point-identification cycle. The two points may be
nonordinary. This includes additional
unibranch singular points of the critical curve.
\end{statement}
\begin{proof}
After P4R.13, the remaining positive ruled case has two
rational normalized conductor components and one graph
cycle. In the stated instance the cycle is provided by
two two-branch points. We first extend the local lower
bound of P4.41 to these points.

At a planar curve point, write its ideal as $(z,h)$.
Its symbolic square equals its ordinary square, so the
rank-one Hessian belongs to $\operatorname{Sym}^2 E_0$,
where $E_0$ is the normalized conormal image. On a branch
put $E_0=\langle t^a e_1,t^b e_2\rangle$ and write the
primitive Hessian direction as $\ell=u e_1+v e_2$.
One of $u,v$ is a unit. Membership of
$t^d\ell^2$ in $\operatorname{Sym}^2E_0$ and projection
onto the quotient by $\ell$ give
\[
 \frac d2\ge\max\{a-\operatorname{ord}u,
                         b-\operatorname{ord}v\},\qquad
 b_{\rm proj}=\min\{a+\operatorname{ord}v,
                         b+\operatorname{ord}u\}.
\]
The sum of the maximum and minimum on the right is
$a+b$. Thus $d/2+b_{\rm proj}\ge a+b$.
The plane different formula proved in P4.34 is
$\sum(a+b)=2\delta+\deg R$. The whole local defect
is consequently at least $3\delta+q$.
A nonordinary planar two-branch point has $\delta\ge2$
and contributes at least six.

For a nonplanar point, the ambient Hessian vanishes: a
nonzero rank-one Hessian would give $f=X^2+g(Y,Z)$ by
the parameter Morse lemma and would put the entire
critical curve in the smooth plane $X=0$.
On a branch of multiplicity $m$, all conormal generators
have order at least $m$, and every Hessian entry has
order at least $m$. Thus $b_{\rm proj}\ge m$, $d\ge m$,
and, since $r=m-1$,
\[
                  d/2+b_{\rm proj}-r\ge1+m/2\ge1.
\]
When $\delta\ge3$, a nonplanar two-branch point
therefore contributes at least five.
The only remaining possibility has $\delta=2$.
The exact excess-one classification in P4.40 identifies
it with a cusp and transverse line,
\[
                  P=(xz,yz,h),\qquad h=y^2-x^3.
\]
Here $P^{(2)}/P^2$ is generated by $zh$ and is
annihilated by the maximal ideal. Write $f=Azh+P^2$.
Along the cusp, the normal quadratic determinant is
$4KG-A^2$, with $K(0)=0$. Its rank is one generically,
so $A(0)=0$ and in fact $f\in P^2$.
The conormal lattices on the two normalized branches
are respectively $\operatorname{diag}(t^2,t^3)$ and
$\operatorname{diag}(t,t)$. Their total colength is seven,
while $\deg R=1$ and $\delta=2$.
Applying the same symmetric-square inequality gives
defect at least $2+7-1=8$.
Every two-branch point thus contributes at least $9/2$,
and equality is possible only at the ordinary node of
P4.41, whose normalization is smooth.

Apply the rank-one Hessian calculation
of P4R.10 to the two normalized projective lines. On each
line the rank-two conormal kernel has degree $2+\deg R$.
The rank-one quotient has Euler characteristic
$3+\deg R-\deg L$. Combining both lines with the
normalization sequence of $C$ gives the exact formula
\[
 \chi(\mathcal O_{C_{\rm sch}})
   =8-\sum_p\left[
       \delta_p+\sum_{i\text{ over }p}
          \left(\frac{d_i}{2}+b_i-r_i\right)+q_p
                 \right].                            \tag{4.42a}
\]
Every extra unibranch singular point has a strictly positive
bracket by the local conormal inequalities of P4R.10.
At a smooth point the bracket is nonnegative. The two
two-branch points contribute at least nine together by the
local calculation just given.

If the horizontal section is not contracted, P4.33 gives
$\chi(\mathcal O_{C_{\rm sch}})=g-1\ge0$, contradicting
(4.42a), whose right side is at most $-1$.
If the section is contracted, that Euler characteristic is
$-1$. Equality must then hold at both two-branch points and every
other bracket must vanish. There are no additional
unibranch singular points. Both cycle points are ordinary
nodes and the normalization is smooth over them by P4.41.
At every smooth point of the
critical curve, P4R.14 gives either the product tacnode
or the displayed $A_1$ normalization degeneration. The
latter has Hessian zero order one and contributes $1/2$,
so it is excluded by the equality in (4.42a). Thus the
normalization is smooth along all of the conductor, and
it is smooth away from the conductor by P4R.6. A contracted
positive-genus horizontal section in its minimal resolution
is impossible. This proves the assertion.
\end{proof}
